% part: intuitionistic-logic
% chapter: soundness-completeness
% section: lindenbaum

\documentclass[../../../include/open-logic-section]{subfiles}

\begin{document}

\olfileid{int}{sc}{lin}

\olsection{Lema Lindenbaum}

Teorema kelengkapan kanggo logika intuisionistik dibuktekake kanthi
nganggep $\Gamma \Proves/ !A$ lan mbangun model kanthi
$\mSat{M}{\Gamma}$ lan~$\mSat/{M}{!A}$.

Ing logika klasik, relasi !!{derivability} bisa direduksi dadi
pangerten konsistensi, amarga !!a{formula}~$!A$ iku !!{derivable}
saka himpunan !!{formula}s yen lan mung yen himpunan iku bebarengan
karo negasi saka~$!A$ ora konsisten. Iki ora bisa ditindakake ing
logika intuisionistik. Ing logika intuisionistik, yen $\lnot!A$ ora
konsisten, kita mung entuk $\Proves \lnot\lnot !A$. Amarga
$\lnot\lnot!A \lif !A$ ora lumaku sacara intuisionistik ing umum,
kita ora bisa nyimpulake yen~$\Proves !A$.

Mula nalika mbangun model~$\mModel{M}$, kita kudu njaga supaya
!!{derivability} kanggo !!{formula}~$!A$ tetep ora kelakon. Dadi kita
ora bisa nggunakake himpunan lengkap $\Gamma^* \supseteq \Gamma$
kanggo mbangun model $\mModel{M}$, amarga ing saben himpunan lengkap
$\Gamma^*$, kita nduweni $\Gamma^* \Proves !A \lor \lnot !A$.

Minangka gantine himpunan lengkap~$\Gamma^*$, kita bakal nggunakake
pangerten himpunan formula kang prim:

\begin{defn}\ollabel{defn:prime}
  Himpunan !!{formula}s~$\Gamma$ iku \emph{prim} yen lan mung yen
  \begin{enumerate}
  \item\ollabel{defn:prime1} $\Gamma$ konsisten, yaiku $\Gamma \Proves/ \lfalse$;
  \item\ollabel{defn:prime2} yen $\Gamma \Proves !A$, mula $!A \in \Gamma$; lan
  \item\ollabel{defn:prime3} yen $!A \lor !B \in \Gamma$, mula $!A \in
    \Gamma$ utawa $!B \in \Gamma$.
  \end{enumerate}
\end{defn}

\begin{lem}[Lema Lindenbaum]
  \ollabel{lem:lindenbaum} Yen $\Gamma \Proves/ !A$, ana
  $\Gamma^* \supseteq \Gamma$ kanthi $\Gamma^*$ prim lan
  $\Gamma^* \Proves/ !A$.
\end{lem}

\begin{proof}
  Upamane $!B_1 \lor !C_1$, $!B_2 \lor !C_2$, \dots, iku enumerasi
  kabeh !!{formula}s kang awujud~$!B \lor !C$. Kita bakal nemtokake
  urutan himpunan !!{formula}s~$\Gamma_n$ kang mundhak, kanthi saben
  $\Gamma_{n+1}$ ditetepake minangka $\Gamma_n$ bebarengan karo siji
  !!{formula} anyar yen isih ana kang kudu ditambahake. $\Gamma^*$
  bakal dadi gabungan kabeh~$\Gamma_n$. !!{formula}s anyar dipilih
  supaya $\Gamma^*$ prim lan tetep $\Gamma^* \Proves/ !A$.
  Tegese, ing saben langkah kita kudu nemokake disjungsi kapisan
  ~$!B_i \lor !C_i$ kanthi:
  \begin{enumerate}
  \item\ollabel{gamma-1} $\Gamma_n \Proves !B_i \lor !C_i$
  \item\ollabel{gamma-2} $!B_i \notin \Gamma_n$ lan $!C_i \notin \Gamma_n$
  \end{enumerate}
  Kita nambahake marang $\Gamma_n$ salah siji $!B_i$ yen
  $\Gamma_n \cup \{!B_i\} \Proves/ !A$, utawa $!C_i$ yen ora.
  Kita kudu nuduhake yen cara iki kasil. Sauntara, tetepa $i(n)$
  minangka $i$ paling cilik kang nyukupi \olref{gamma-1}
  lan~\olref{gamma-2}.
  
  Tetepa $\Gamma_0 = \Gamma$ lan
  \[
  \Gamma_{n+1} = \begin{cases}
    \Gamma_n \cup \{!B_{i(n)}\} &
    \text{yen $\Gamma_n \cup \{!B_{i(n)}\} \Proves/ !A$} \\
    \Gamma_n \cup \{!C_{i(n)}\} & \text{yen ora}
    \end{cases}
  \]
  Yen $i(n)$ ora ditetepake, yaiku saben $\Gamma_n \Proves !B \lor !C$,
  salah siji $!B \in \Gamma_n$ utawa $!C \in \Gamma_n$, kita nemtokake
  $\Gamma_{n+1} = \Gamma_n$. Saiki tetepa
  $\Gamma^* = \bigcup_{n=0}^\infty \Gamma_n$.

  Sepisan, kita nuduhake yen kanggo kabeh $n$, $\Gamma_n \Proves/ !A$.
  Kita nganggo induksi tumrap~$n$. Kanggo $n = 0$, pratelan iki
  bener miturut hipotesis lema, yaiku $\Gamma \Proves/ !A$.
  Kanggo $n \ge 0$, kita kudu nuduhake yen $\Gamma_n \Proves/ !A$,
  mula $\Gamma_{n+1} \Proves/ !A$. Yen $i(n)$ ora ditetepake,
  $\Gamma_{n+1} = \Gamma_n$ lan ora ana kang kudu dibuktekake.
  Mula upamane $i(n)$ ditetepake. Supaya prasaja, tetepa $i = i(n)$.
  
  Kita bakal mbuktekake kontrapositif pratelan iki. Upamane
  $\Gamma_{n+1} \Proves !A$. Miturut konstruksi,
  $\Gamma_{n+1} = \Gamma_n \cup \{!B_i\}$ yen
  $\Gamma_n \cup \{!B_i\} \Proves/ !A$, utawa yen ora,
  $\Gamma_{n+1} = \Gamma_n \cup \{!C_i\}$. Cetha ora bisa kang
  kapisan, amarga banjur $\Gamma_{n+1} \Proves/ !A$. Mula
  $\Gamma_n \cup \{!B_i\} \Proves !A$ lan
  $\Gamma_{n+1} = \Gamma_n \cup \{!C_i\}$. Miturut definisi $i(n)$,
  kita nduweni $\Gamma _n \Proves !B_i \lor !C_i$.
  Kita nduweni $\Gamma_n \cup \{!B_i\} \Proves !A$. Uga
  $\Gamma_{n+1} = \Gamma_n \cup \{!C_i\} \Proves !A$.
  Mula $\Gamma_n \Proves !A$, kaya kang arep dituduhake.

  Yen $\Gamma^* \Proves !A$, ana subhimpunan winates
  $\Gamma' \subseteq \Gamma^*$ kanthi $\Gamma' \Proves !A$.
  Saben $!D \in \Gamma'$ kudu ana ing $\Gamma_i$ kanggo sawenehing~$i$.
  Pilihen $n$ kang ora luwih cilik tinimbang kabeh indeks iku;
  yen ora ana indeks amarga himpunane kosong, pilih nol.
  Amarga $\Gamma_i \subseteq \Gamma_n$ yen $i \le n$,
  $\Gamma' \subseteq \Gamma_n$. Nanging banjur $\Gamma_n \Proves !A$,
  nalisir bukti ing ndhuwur yen $\Gamma_n \Proves/ !A$.

  Pungkasan, kita nuduhake yen $\Gamma^*$ prim, yaiku nyukupi syarat
  \olref{defn:prime1}, \olref{defn:prime2}, lan~\olref{defn:prime3}
  saka \olref{defn:prime}.

  Sepisan, $\Gamma^* \Proves/ !A$, mula $\Gamma^*$ konsisten.
  Dadi \olref{defn:prime1} kasembadan.

  Saiki kita nuduhake yen $\Gamma^* \Proves !B \lor !C$, mula salah
  siji $!B \in \Gamma^*$ utawa $!C \in \Gamma^*$. Iki mbuktekake
  \olref{defn:prime3}, amarga yen $!B \lor !C \in \Gamma^*$,
  uga $\Gamma^* \Proves !B \lor !C$. Mula upamane
  $\Gamma^* \Proves !B \lor !C$ nanging $!B \notin \Gamma^*$ lan
  $!C \notin \Gamma^*$. Amarga $\Gamma^* \Proves !B \lor !C$,
  $\Gamma_n \Proves !B \lor !C$ kanggo sawenehing~$n$.
  $!B \lor !C$ ana ing enumerasi kabeh disjungsi, upamane minangka
  $!B_j \lor !C_j$. $!B_j \lor !C_j$ nyukupi sipat-sipat ing definisi
  $i(n)$, yaiku $\Gamma_n \Proves !B_j \lor !C_j$, dene
  $!B_j \notin \Gamma_n$ lan $!C_j \notin \Gamma_n$.
  Disjungsi iki tetep nyukupi syarat ing saben tahap sabanjure nganti
  dipilih. Ing saben tahap sadurunge dipilih, disjungsi $!B_i \lor !C_i$
  kanthi indeks luwih cilik dipilih. Indeks kang luwih cilik mung ana
  winates, lan saben disjungsi iku mung bisa dipilih sapisan, amarga
  saben dipilih kita nambahake salah siji $!B_i$ utawa $!C_i$.
  Mula ing sawenehing tahap~$m$, kita bakal nduweni $j = i(m)$.
  Nanging banjur salah siji $!B \in \Gamma_{m+1}$ utawa
  $!C \in \Gamma_{m+1}$, nalisir asumsi yen $!B \notin \Gamma^*$
  lan $!C \notin \Gamma^*$.

  Saiki upamane $\Gamma^* \Proves !B$. Mula $\Gamma^* \Proves !B \lor !B$.
  Nanging kita lagi wae mbuktekake yen $\Gamma^* \Proves !B \lor !B$,
  mula $!B \in \Gamma^*$. Dadi $\Gamma^*$ nyukupi
  ~\olref{defn:prime2} saka~\olref{defn:prime}.
\end{proof}

\begin{prob}
Tuduhna yen $\Gamma \Proves/ \bot$, mula $\Gamma$ konsisten ing logika
klasik, yaiku ana valuasi kang ndadekake kabeh formula ing~$\Gamma$ bener.
\end{prob}

\end{document}
